\documentclass[exercice]{thedocument}%
\usepackage{texmaths-tikz}%
\startdatakeys
\source{}
\cycle{}
\competencesDuSocle{}
\connaissancesRequises{}
\competencesTravaillees{}
\enddatakeys
\begin{document}%
    Les droites (d$_1$) et (d$_2$) sont parallèles.\par\noindent
    Donner toutes les mesures d'angles que l'on peut trouver en justifiant
    les réponses. (utiliser des couleurs pour marquer les angles de même mesure).
    \begin{center}
        \begin{tikzpicture}
            \coordinate(A)at(0,0);
            \coordinate(A0)at(-1,0);
            \coordinate[label=below right:(d$_2$)](A1)at(3,0);
            \coordinate(B)at(2,2);
            \coordinate[label=above right:(d$_1$)](B0)at(-0.5,2);
            \coordinate(B1)at(3.5,2);
            \path(B)--(A)coordinate[pos=1.2](C);
            \tkzFillAngle[fill=\currentcolor,size=0.6](A0,A,C)
            \tkzLabelAngle(A0,A,C){55$^\circ$}
            \tkzDrawLine[add=0.6 and 0.6,line width=1pt](A,B)
            \tkzDrawLines[line width=1pt](A0,A1 B0,B1)
        \end{tikzpicture}
    \end{center}
\end{document}
